%%%%%%%%%%  scan idx 429  |  folio 416  %%%%%%%%%%
% Exposé VII: back to the volume's Sym and Proj
\renewcommand{\Sym}{\operatorname{Sym}}
\renewcommand{\Projz}{\operatorname{Proj}_0}
\section*{Exposé VII. Immersions régulières et calcul du \texorpdfstring{$K^{\bullet}$}{K0} d'un schéma éclaté}
\begin{center}
par Pierre \textsc{Berthelot}
\end{center}

\section*{1. Généralités sur les immersions régulières}

La notion classique d'immersion régulière (EGA IV 16.9.2) s'avérant inadéquate dans le cas non noethérien, nous allons définir une notion plus faible, mais donnant lieu aux propriétés souhaitées, et qui équivaut à la définition usuelle dans le cas noethérien.

\textbf{Définition 1.1.} Soient $A$ un anneau commutatif, $E$ un $A$-module projectif de type fini, $u:E\longrightarrow A$ un homomorphisme $A$-linéaire. On dit que $u$ est régulier si le complexe de Koszul $K_\bullet(u)$ défini par $u$ est acyclique en dimension $>0$, donc est une résolution de $B=A/u(E)$.

Rappelons que le complexe de Koszul associé à $u$ est défini comme suit: le $A$-module gradué $K_\bullet(u)$ est l'algèbre extérieure $\Lambda(E)$, et l'opérateur bord la multiplication intérieure par $u$ considéré comme élément du dual $\check E$ : c'est une antidérivation de degré $-1$ de $\Lambda(E)$, égale à $u$ sur $E$.

\textbf{Lemme 1.2.} Soit $u:E\longrightarrow A$ un homomorphisme régulier. Soient $A'$ une $A$-algèbre, $u':E'\longrightarrow A'$ l'homomorphisme déduit de $u$ par changement de base; on suppose $u'$ régulier. On pose $J=u(E)$, $J'=JA'=u(E')$. Alors on a les relations suivantes pour tout entier $n\geq0$:
\begin{align}
 J^n/J^{n+1}\otimes_AA'&\xrightarrow{\sim}J'^n/J'^{n+1}, \tag{1.2.1}\\
 \Tor_i^A(J^n/J^{n+1},A')&=0\quad\text{pour }i\geq1; \tag{1.2.2}\\
 J^n\otimes_AA'&\xrightarrow{\sim}J'^n, \tag{1.2.3}\\
 \Tor_i^A(J^n,A')&=0\quad\text{pour }i\geq1. \tag{1.2.4}
\end{align}

a) On suppose d'abord que les $J^n/J^{n+1}$ sont projectifs sur $B=A/J$; comme $J$ est de type fini, ils sont localement libres de type fini sur $B$.
